\documentclass[../../../include/open-logic-section]{subfiles}

\begin{document}

\olfileid{sth}{z}{power}
\olsection{ఘాత సమితులు}

మరో స్వీకృతంతో ముందుకు
వెళ్దాం:

\begin{axiom}[ఘాత సమితులు]
ఏ సమితి $A$కైనా,
$\Pow{A} = \Setabs{x}{x \subseteq A}$ సమితి ఉంటుంది.
\[
	\forall A \exists P \forall x(x \in P \liff (\forall z \in x)z \in A)
\]
\end{axiom}

దీనికి మన సమర్థన చాలా
సరళం. $A$ సమితి $S$
దశలో ఏర్పడిందనుకుందాం.
అప్పుడు $A$లోని మూలకాలన్నీ
$S$కు ముందే అందుబాటులో
ఉన్నాయి (\stagesacc\ ప్రకారం).
అందువల్ల వేరుచేయడం
సమర్థనలో మాదిరిగానే,
$A$లోని ప్రతి ఉపసమితీ
$S$ దశ నాటికి ఏర్పడుతుంది.
కాబట్టి అవన్నీ $S$
తరువాతి ఏ దశలోనైనా
ఒకే సమితిగా ఏర్పడడానికి
అందుబాటులో ఉంటాయి.
$S$ చివరి దశ కాదు
(\stagessucc\ ప్రకారం);
అందుకే అటువంటి దశ
ఒకటి ఉందని తెలుసు.
కాబట్టి $\Pow{A}$ ఉంటుంది.

ఘాత సమితుల స్వీకృతానికి
ఒక చక్కని పర్యవసానం ఇది:

\begin{prop}\label{thm:Products}
ఏ సమితులు $A, B$ ఇచ్చినా,
వాటి కార్టీజియన్ లబ్ధం
$A \times B$ ఉంటుంది.
\end{prop}

\begin{proof}
ఘాత సమితుల స్వీకృతం,
\olref[sth][z][pairs]{prop:pairsconsequences}
ప్రకారం $\Pow{\Pow{A \cup B}}$
సమితి ఉంటుంది. కాబట్టి
వేరుచేయడం ప్రకారం ఈ
సమితి ఉంటుంది:
\[
	C = \Setabs{z \in \Pow{\Pow{A \cup B}}}{(\exists x \in A)(\exists y \in B) z = \tuple{x, y}}.
\]
ఇప్పుడు ఏ $x \in A$,
$y \in B$కైనా
\olref[sth][z][pairs]{prop:pairsconsequences}
ప్రకారం $\tuple{x, y}$
ఉంటుంది. ఇంకా, $x, y
\in A \cup B$ కాబట్టి,
$\{x\}, \{x, y\} \in \Pow{A \cup B}$;
మరియు $\tuple{x,y} \in \Pow{\Pow{A \cup B}}$.
కాబట్టి $A \times B = C$.
\end{proof}

ఈ నిరూపణలో ఘాత సమితుల
స్వీకృతం వేరుచేయడంతో
కలిసి పనిచేస్తుంది.
ఇది ఆశ్చర్యం కాదు.
వేరుచేయడం లేకపోతే,
ఘాత సమితుల స్వీకృతం
అంత \emph{శక్తివంతమైన}
సూత్రం కాదు. ఏ సమితికి
ఏ ఉపసమితులు ఉన్నాయో
వేరుచేయడం చెబుతుంది;
అందువల్ల ప్రతి ఘాత
సమితి ఎంత ``నిండుగా''
ఉంటుందో నిర్ణయిస్తుంది.

\begin{prob}
ఏ సమితులు $A, B$కైనా:
(i) నిర్వచన సమితి $A$,
విలువల పరిధి $B$ గల
అన్ని సంబంధాల సమితి
ఉంటుందని; (ii) $A$
నుంచి $B$కు గల
అన్ని ప్రమేయాల సమితి
ఉంటుందని చూపండి.
\end{prob}

\begin{prob}
$A$ ఒక సమితి, $\sim$
దానిపై సమానతా సంబంధం
అనుకుందాం. $A$పై $\sim$
ప్రకారం $A$లోని సమానతా వర్గాల
సమితి, అంటే
$\equivclass{A}{\sim}$,
ఉంటుందని నిరూపించండి.
\end{prob}

\end{document}
